% !PRINTED 667
%% [running head, not body text -- commented out 2026-08-07] Page furniture (running
%% head, folio, signature) is excluded from this edition by standing policy.
%% \S.\,184. \hfill \gesperrt{Fall eines geraden $m$.} \hfill 667
\weberpage{667}
%<<CONT
Wenn wir also nach (5)

\[
\mathfrak{E}^x \mathfrak{E}(r) = e(r)
\]

setzen, so folgt

\weq{7}{\mathfrak{E}(r) = e(r)^m.}

Durch (7) ist aber der Satz 2. bewiesen und damit zugleich f\"ur ein
\gesperrt{ungerades $m$ das ganze Theorem I.}

Auch dieser Schluss versagt f\"ur ein gerades $m$, weil dann die
Gleichung (6) nicht mehr l\"osbar ist.

\begin{center}
\S.\,184.

\gesperrt{Vorl\"aufiges \"uber den Fall eines geraden $m$.}
\end{center}

F\"ur den Fall, dass $m$ eine Potenz von 2 ist, schlagen wir einen
ganz anderen Weg ein, \"uber den hier zun\"achst einige vorl\"aufige
Bemerkungen Platz finden m\"ogen.

Wir haben schon im \S.\,182 gesehen, dass das erste Lemma bewiesen
ist, wenn wir nachweisen k\"onnen, dass irgend eine Potenz von
$\varphi$, deren Exponent nicht durch $q$ theilbar ist, zur
Hauptclasse geh\"ort. Dabei ist von den beiden Eigenschaften des
Functionals $\varphi$ nur die erste benutzt, dass die $m^{\text{te}}$
Potenz von $\varphi$ in der Hauptclasse enthalten ist.

Nun kennen wir nach den allgemeinen S\"atzen in \S.\,154 in jedem
K\"orper einen Exponenten $h$, n\"amlich die Anzahl der
Idealclassen, f\"ur den $\varphi^h$ zur Hauptclasse geh\"ort, wenn
$\varphi$ ein beliebiges Functional in diesem K\"orper ist.

Wenn wir also beweisen k\"onnen, dass die \gesperrt{Classenzahl $h$ im K\"orper
$\Omega_m$, wenn $m$ eine Potenz von 2 ist, eine ungerade Zahl ist,}
so ist damit ohne alles Weitere das Lemma 1. bewiesen.

Dies soll das Ziel unserer Betrachtungen in den n\"achsten
Abschnitten sein.

In Bezug auf das zweite Lemma liegt die Sache \"ahnlich. Hier kann
man aus den Voraussetzungen in \S.\,183, 2. ganz wie oben die Formel
\S.\,183, (5) herleiten, aus der aber nur zu schliessen ist, dass
$\mathfrak{E}(r)$ die $\frac{1}{2} m^{\text{te}}$ Potenz einer
reellen Einheit $e(r)$ ist.

Nehmen wir also das erste Lemma als bewiesen an, so ergiebt die
Formel \S.\,181, (14):

\weq{1}{(r^n, x_0) = \varrho_n \sqrt{e(r^n)}\ \alpha_n\ \Pi(r^{m_2 n}, \eta),}
%CONT>>
